\section{Entrywise divergence and the Bures value}\label{sec:bures}
In this section $k=nd$, and $q=(Q_1,\ldots,Q_d)$ is the horizontal
concatenation of $n\times n$ matrix-valued measure blocks. We use
$D^*\Phi=(\partial_1\Phi,\ldots,\partial_d\Phi)$.
The continuity equation remains the weak formulation \eqref{eq:ce}.
For smooth admissible densities with $F=(F_1,\ldots,F_d)$, the block
divergence and continuity equation are
\[
 (\operatorname{Div}F)_{ij}
   =\sum_{a=1}^d\partial_a(F_a)_{ij},\qquad
 \partial_tM+\sum_{a=1}^d\partial_a\sym F_a=\sym Z.
\]
The second identity is the density form of \eqref{eq:ce} in the interior.
For an $n\times d$ matrix field $\Phi$, row divergence instead means
$(\operatorname{div}\Phi)_i=\sum_{a=1}^d\partial_a\Phi_{ia}$;
it is the test-side operator used in the Kantorovich--Bures examples.
Every spatially constant symmetric test is annihilated by $D^*$, a property
stronger than $D^*I_n=0$. Fix $H\in\Spp^n$.

For PSD matrices define the usual Bures--Wasserstein expression
\[
 d_B^2(X,Y)=\tr X+\tr Y-2\tr\sqrt{X^{1/2}YX^{1/2}}.
\]
This expression and the orthogonal Procrustes characterization are standard
in matrix geometry \cite{BhatiaJainLim}. Weighted reaction congruence is already
present in \cite[Remark 3.11]{LiZou} and is related to \cite{HanEtAl}.

\begin{lemma}[Pure reaction with PSD endpoints]\label{lem:pure}
Let $\sigma$ be a probability Radon measure on $\Omega$ and
$A_0,A_1\succeq0$. For endpoints $A_0\sigma,A_1\sigma$, the minimum
of the original action over admissible paths with $q=0$ equals
\begin{equation}\label{eq:buresvalue}
 E_H(A_0,A_1)=2d_B^2(H^{-1}A_0H^{-1},H^{-1}A_1H^{-1}).
\end{equation}
The matrices $A_i$ may be singular, and the minimum is attained.
\end{lemma}
\begin{proof}
Write $X_i=H^{-1}A_iH^{-1}$. Choose factors $P_iP_i^T=X_i$ with
$P_0=X_0^{1/2}$ and $\norm{P_1-P_0}_F^2=d_B^2(X_0,X_1)$.
Here is an explicit construction covering singular endpoints. For $\delta>0$ put
\[
 S_\delta=(X_0+\delta I)^{1/2},\qquad
 K_\delta=(S_\delta X_1S_\delta)^{1/2},\qquad
 P_\delta=S_\delta^{-1}K_\delta.
\]
Then $P_\delta P_\delta^T=S_\delta^{-1}K_\delta^2S_\delta^{-1}=X_1$
and $\tr(S_\delta P_\delta)=\tr K_\delta$. In particular
$\norm{P_\delta}_F^2=\tr X_1$, so finite-dimensional compactness supplies
a convergent subsequence as $\delta\downarrow0$; denote its limit by $P_1$.
The PSD square-root map is continuous: for a convergent sequence of PSD
matrices its square roots are bounded, and every subsequential limit is PSD
and squares to the limiting matrix, whose PSD square root is unique.
Consequently the factor and trace identities pass to the limit, giving
\[
 P_1P_1^T=X_1,\qquad
 \tr(P_0P_1)=\tr\bigl((X_0^{1/2}X_1X_0^{1/2})^{1/2}\bigr).
\]
Expanding
$\norm{P_1-P_0}_F^2=\tr X_1+\tr X_0-2\tr(P_0P_1)$ proves the required
identity and the nonnegativity of $d_B^2$.

Put $\Delta=P_1-P_0$, $P_s=P_0+s\Delta$, and
\[
 M_s=HP_sP_s^TH,\qquad Z_s=2HP_s\Delta^TH.
\]
Differentiation gives
$\dot M_s=H(\Delta P_s^T+P_s\Delta^T)H=\sym Z_s$. The path
$(M_s\,ds\otimes\sigma,0,Z_s\,ds\otimes\sigma)$ has endpoints
$A_i\sigma$ and satisfies the original weak equation. Symmetric tests pair
with $Z_s$ through $\sym Z_s$. For interior $x$, the product rule and the
fundamental theorem of calculus in time applied to $\Phi(s,x)\cdot M_s$ give the endpoint
identity. For $x\in\partial\Omega$, choose interior points $x_m\to x$,
possible because $\Omega$ is the closure of its interior. Uniform continuity
of $\Phi$ and $\partial_s\Phi$ on compact $Q$ gives uniform convergence
in $s$ along these points. Since $M_s$ and $\dot M_s$ are bounded, the
interior identity passes to the limit. Integrating this identity in the
finite measure $\sigma$ proves \eqref{eq:ce} also for boundary-supported
$\sigma$.
For an admissible dual triple $(A,B,C)$, the nonnegative transport Gram
term may be dropped after pairing its constraint with $M_s$. With
$Y=(CH)^THP_s$, cyclicity of trace gives
$\tr(CH^2C^TM_s)=\norm Y_F^2$ and $C\cdot Z_s=2\ip Y\Delta$.
Thus
\[
 A\cdot M_s+C\cdot Z_s
 \le-\tfrac12\norm Y_F^2+2\ip Y\Delta
 =2\norm\Delta_F^2-\tfrac12\norm{Y-2\Delta}_F^2
 \le2\norm\Delta_F^2.
\]
Integrating against $ds\otimes\sigma$ and taking the continuous-dual
supremum in \eqref{eq:continuousdual} proves that this path has
pure-reaction action at most $2d_B^2(X_0,X_1)$. The estimate applies even
where $M_s$ is singular.

For the matching lower bound, regularize both endpoints and set
\[
 S_\delta=(X_0+\delta I)^{1/2},\qquad
 T_\delta=S_\delta^{-1}
   [S_\delta(X_1+\delta I)S_\delta]^{1/2}S_\delta^{-1}.
\]
The symmetric matrix $T_\delta$ is positive definite. Thus
$N_s=(1-s)I+sT_\delta$ is invertible for $s\in[0,1]$, and
\[
 \Psi_s=2(T_\delta-I)N_s^{-1},\qquad
 \Phi_s=H^{-1}\Psi_sH^{-1}
\]
are time-dependent symmetric tests. Indeed, $N_s$ is a polynomial in
$T_\delta$, so $N_s^{-1}$ commutes with $T_\delta-I$. Differentiating
$N_sN_s^{-1}=I$ yields
\[
 \dot\Psi_s=-2(T_\delta-I)N_s^{-1}(T_\delta-I)N_s^{-1}
            =-\tfrac12\Psi_s^2.
\]
Since $H^{-1}H^2H^{-1}=I$, congruence by $H^{-1}$ gives
$\dot\Phi_s+\Phi_sH^2\Phi_s/2=0$. Hence
$(\dot\Phi_s,0,\Phi_s)$ satisfies the dual constraint. The original weak
equation and the defining support inequality bound every pure-reaction
competitor below by $\Phi_1\cdot A_1-\Phi_0\cdot A_0$; the spatial
integrals reduce to these pairings because $\sigma(\Omega)=1$.
Write $\widehat K_\delta=[S_\delta(X_1+\delta I)S_\delta]^{1/2}$ and
$c_\delta=\tr\widehat K_\delta$. Direct multiplication gives
\[
 T_\delta S_\delta^2T_\delta
 =S_\delta^{-1}\widehat K_\delta^2S_\delta^{-1}=X_1+\delta I,
 \qquad \tr(T_\delta S_\delta^2)=\tr\widehat K_\delta=c_\delta.
\]
Also $T_\delta^{-1}(X_1+\delta I)=S_\delta^2T_\delta$, whose trace is
$c_\delta$. Subtracting the regularization terms therefore gives
\[
 \tr(T_\delta X_0)=c_\delta-\delta\tr T_\delta,
 \qquad
 \tr(T_\delta^{-1}X_1)=c_\delta-\delta\tr T_\delta^{-1}.
\]
Since $\Psi_0=2(T_\delta-I)$ and
$\Psi_1=2(I-T_\delta^{-1})$, the endpoint pairing is the sum of
\[
 2\{\tr X_0+\tr X_1-
   2\tr\bigl([S_\delta(X_1+\delta I)S_\delta]^{1/2}\bigr)\}
\]
and $2\delta\tr(T_\delta+T_\delta^{-1})\ge0$.
Drop this nonnegative second term before passing to the limit. The first
term converges to $2d_B^2(X_0,X_1)$ by continuity of the PSD square root,
so every competitor has at least that action. This proves the lower bound,
including singular endpoints. The same space-constant certificate gives the lower bound for
any original transport problem whose operator annihilates all spatially
constant symmetric fields.
\end{proof}

\paragraph{A noncommuting inflation example.}
We compare the inflation expression in
\cite[Proposition 4.4, equation (4.12)]{LiZou} with the variational value
for the following noncommuting pair. Take
$x_0\in\operatorname{int}\Omega$, $e=(1,0)^T$, $v=(1,1)^T/\sqrt2$,
$P_0=ee^T$, $P_1=vv^T$, and $H=I_2$.
For $\mathsf G=I_2\delta_{x_0}$, the inflated endpoints
$\mathsf G^{P_i}:=P_i\mathsf G P_i$ equal $P_i\delta_{x_0}$.
For these endpoints, that expression equals $2\tr((P_1-P_0)^2)=2$.
Put $w=v-e$ and $u_s=e+sw$. Since
$|u_s|^2=1-(2-\sqrt2)s(1-s)>0$, the zero-flux path with mass
$u_su_s^T\,ds\otimes\delta_{x_0}$ and reaction
$2u_sw^T\,ds\otimes\delta_{x_0}$ is admissible:
$\partial_s(u_su_s^T)=\sym(2u_sw^T)$.
Its local support cost is
\[
 j_{L,I_2}(u_su_s^T,0,2u_sw^T)
 =\sup_C\left\{2w\cdot C^Tu_s-\tfrac12|C^Tu_s|^2\right\}
 =2|w|^2=4-2\sqrt2<2.
\]
Indeed, the maximizing choice of the mass dual variable saturates its
constraint, the transport dual variable can be set to zero, and
$C^Tu_s$ ranges over $\R^2$.
For entrywise divergence, Lemma~\ref{lem:pure} and its space-constant
lower certificate give the exact minimum
$2d_B^2(P_0,P_1)=4-2\sqrt2<2$.

\begin{theorem}[Radon hub recovery]\label{thm:hub}
Suppose $\Omega$ is convex, choose $x_*\in\operatorname{int}\Omega$, and
put $A_i=G_i(\Omega)$,
\[
 M_i=\int_\Omega|x-x_*|^2\,d\tr G_i(x),\qquad
 B=(\sqrt{M_0}+\sqrt{M_1})/\sqrt2.
\]
For arbitrary finite PSD Radon endpoints and every $t>0$,
\begin{equation}\label{eq:hub}
 E_H(A_0,A_1)\le V(tI_{nd},H)
 \le\bigl(\sqrt{E_H(A_0,A_1)}+B/t\bigr)^2.
\end{equation}
Consequently $W(I_{nd},H)=E_H(A_0,A_1)$. For SPD $L$ with
$L^2\succeq m^2I_{nd}$, $m>0$, the same bound holds with $t$ replaced by $m$.
\end{theorem}
\begin{proof}
Throughout this proof abbreviate $E_H:=E_H(A_0,A_1)$.
The space-constant certificates in Lemma~\ref{lem:pure} are admissible
for entrywise divergence. They bound every original competitor below by
$E_H(A_0,A_1)$, since their endpoint pairings depend only on the total
matrices $A_i$. This proves the lower bound.

For the upper bound, we contract the endpoints to a common hub and
connect their total matrices there by pure reaction. For a single endpoint
$G$, take for $\lambda$ the sum of its coordinate total variations and
write $G=N\lambda$. Since
$G^{ab}=G^{ba}$, uniqueness of Radon--Nikodym derivatives gives
$N=N^T$ almost everywhere. For each vector $v$ in a countable dense subset
of $\R^n$, positivity of the scalar measure $v^TGv$ implies
$v^TNv\ge0$ almost everywhere. Removing the union of these countably many
null sets and using continuity in $v$ gives $N\succeq0$ almost everywhere.
Define $F_s(x)=(1-s)x+sx_*$, $v(x)=x_*-x$, and
$\mathcal F(s,x)=(s,F_s(x))$. The contracted measures are exactly
\[
 \widehat G=\mathcal F_\#(N\,ds\,d\lambda),\qquad
 \widehat Q_\ell=\mathcal F_\#(v_\ell N\,ds\,d\lambda),\qquad
 \widehat R=0.
\]
Their coordinates are finite signed Radon measures: $v$ is bounded on
compact $\Omega$, every coordinate of $N$ is integrable, and continuous
pushforward preserves finiteness and the Radon property on compact spaces.
For a closed-domain test $\Phi$, the time derivative of
$\Phi(s,F_s(x))$ is
$\partial_s\Phi+\sum_\ell v_\ell\partial_\ell\Phi$ at $(s,F_s(x))$.
To justify this also when $x\in\partial\Omega$, choose
$B(x_*,r)\subset\operatorname{int}\Omega$. Convexity gives
$B(F_s(x),sr)\subset\Omega$ for $s>0$, so the ordinary chain rule applies
there; the continuous extensions of the derivatives allow application of the
fundamental theorem of calculus down to $s=0$. Pairing with $N$ and integrating
therefore gives
\[
 \int_Q\partial_s\Phi\cdot d\widehat G+
          \sum_\ell\partial_\ell\Phi\cdot d\widehat Q_\ell
 =\int_\Omega[\Phi(1,x_*)-\Phi(0,x)]\cdot N(x)\,d\lambda(x).
\]
Fubini is applicable because the derivatives and $v$ are bounded and
$\tr N$ is integrable. This is precisely \eqref{eq:ce} from $G$ to
$G(\Omega)\delta_{x_*}$, including boundary atoms.

To bound the pre-pushforward cost, factor $N=SS^T$. For an admissible
dual triple $(A,(B_\ell)_\ell,C)$, dropping the reaction Gram term and
completing the square give
\[
 A\cdot N+\sum_\ell B_\ell\cdot(v_\ell N)
 \le\sum_\ell\left(-\frac{t^2}{2}\norm{B_\ell^TS}_F^2
       +v_\ell\ip{S}{B_\ell^TS}\right)
 \le\frac{|v|^2\tr N}{2t^2}.
\]
Here $\ip{S}{B_\ell^TS}=B_\ell\cdot N$ and
$\norm S_F^2=\tr N$, so the second inequality is the scalar-vector square
completion with maximizing value $v_\ell^2\norm S_F^2/(2t^2)$.
Lemma~\ref{lem:scaling} permits integration against $ds\otimes\lambda$.
For completeness, if $(A,B,C)$ is any continuous admissible field on the
image space, its composition with $\mathcal F$ is a continuous admissible
field on the source: the dual constraint depends only on its values and
on the fixed weights. The signed pushforward identity gives, for example,
\[
 \int_Q A\cdot d\widehat G+\sum_\ell B_\ell\cdot d\widehat Q_\ell
 =\int_Q\left[(A\circ\mathcal F)\cdot N+
      \sum_\ell(B_\ell\circ\mathcal F)\cdot v_\ell N\right]ds\,d\lambda.
\]
The preceding pointwise bound and \eqref{eq:continuousdual} show that the
contraction cost is at most $M_G/(2t^2)$, where
$M_G=\int|x-x_*|^2d\tr G$. This argument also covers collisions of distinct
source points at the hub.
For expansion, run the curve backwards and reverse the flux. The cost is
unchanged, since the dual constraint is invariant under $B\mapsto-B$.

We use the time-rescaling rule of \cite[Lemma 3.15]{LiZou} to concatenate
these paths, with its measure-level form written explicitly below.
For a unit-time piece $\mu=(G,q,R)$ and an interval $[a,a+r]$ of length
$r>0$, let $\Theta(s,x)=(a+rs,x)$ and set
\[
 G'=r\Theta_\#G,\qquad q'=\Theta_\#q,\qquad R'=\Theta_\#R.
\]
For a closed-domain test $\Phi$, $\Psi=\Phi\circ\Theta$ satisfies
$\partial_s\Psi=r(\partial_t\Phi)\circ\Theta$ and
$D^*\Psi=(D^*\Phi)\circ\Theta$. Substitution in the original equation
proves the compressed weak equation with the same spatial endpoints at
times $a,a+r$. If $(g,z,\zeta)$ are densities of $\mu$ with respect to
$\rho$, the densities of the compressed path with respect to
$\Theta_\#\rho$ are $(rg,z,\zeta)\circ\Theta^{-1}$. Reference-measure
independence and \eqref{eq:quadratic-scaling} give the exact identity
\[
 \J_{tI,H}(G',q',R')
 =\int j_{tI,H}(rg,z,\zeta)\,d\rho
 =r^{-1}\J_{tI,H}(G,q,R).
\]

Apply this operation to the first contraction, the pure-reaction path at
the hub with $\sigma=\delta_{x_*}$, and the second expansion, using
consecutive intervals of lengths $r_0,r_1,r_2>0$ summing to one. Each piece
has a Lebesgue factor in time and hence assigns zero to its endpoint slices.
Their coordinate sums therefore have no duplicated junction mass. Summing
the three weak equations cancels their two internal endpoint pairings and
gives \eqref{eq:ce}. Off the junction slices the measures have disjoint
time supports, so reference-measure independence makes the action of the
sum the sum of the piece actions. Consequently the action is at most
\[
 \frac{M_0}{2t^2r_0}+\frac{E_H}{r_1}+\frac{M_1}{2t^2r_2}.
\]
For the three nonnegative costs
$(c_0,c_1,c_2)=(M_0/(2t^2),E_H,M_1/(2t^2))$, Cauchy--Schwarz gives
$(\sum_i\sqrt{c_i})^2\le(\sum_i c_i/r_i)(\sum_i r_i)$.
Thus the infimum over positive durations summing to one is
$(\sqrt{E_H}+B/t)^2$. If all three costs are positive, it is attained by
durations proportional to their square roots; if some, but not all, vanish,
it is approached by positive durations whose zero-cost components tend to
zero; if all three vanish, every such choice gives zero. Taking infima gives the upper bound. Along $t\to\infty$, the two bounds
force $V(tI_{nd},H)\to E_H$. The fixed-ray limit in
Theorem~\ref{thm:ray} identifies the same limit as $W(I_{nd},H)$.
For the general SPD bound, $LL^T\succeq m^2I_{nd}$ and the unchanged
reaction Gram imply $\J_{L,H}\le\J_{mI,H}$ by the support-set
comparison with $\alpha=1$.
\end{proof}
